\section*{Aktivitas Pembelajaran dan Praktikum}
\addcontentsline{toc}{section}{Aktivitas Pembelajaran dan Praktikum}

\begin{table}[H]
\centering
\caption{Rencana aktivitas Bab VIII}
\label{tab:ch08-aktivitas}
\begin{tabularx}{\textwidth}{@{}llX@{}}
\toprule
\textbf{Sesi} & \textbf{Bentuk} & \textbf{Kegiatan} \\
\midrule
1 & Ceramah & CFG, BNF, dan derivasi \\
2 & Latihan & Menurunkan string dan menggambar pohon urai \\
3 & Lokakarya & Konversi EBNF ke BNF \\
4 & Praktikum & P2-A: tata bahasa Subset-C dan peer review \\
\bottomrule
\end{tabularx}
\end{table}

\subsection*{P2-A: Tata Bahasa Subset-C}
\textbf{Tujuan.} Menulis tata bahasa EBNF untuk Subset-C, lalu
mengonversinya ke BNF.

\textbf{Langkah kerja.}
\begin{enumerate}[label=\arabic*.,leftmargin=*]
  \item Lengkapi \file{spec/grammar.ebnf} dengan aturan untuk seluruh
        konstruksi yang didukung.
  \item Turunkan lima string contoh: ekspresi, deklarasi, if, while,
        dan fungsi.
  \item Konversi EBNF ke BNF pada \file{spec/grammar.bnf}.
  \item Gambar pohon urai untuk dua string.
  \item Tukar berkas dengan teman untuk peer review.
\end{enumerate}

\textbf{Artefak.} Berkas \file{spec/grammar.ebnf},
\file{spec/grammar.bnf}, lima derivasi, dua pohon urai, dan catatan
hasil uji.

\textbf{Definisi selesai.}
\begin{checklist}
  \item EBNF mencakup semua konstruksi: deklarasi, ekspresi, dan
        pernyataan.
  \item Lima derivasi benar dan terdokumentasi.
  \item BNF konsisten dengan EBNF.
  \item Dua pohon urai benar dan sesuai derivasi.
  \item Berkas grammar dan bukti pengujian dikumpulkan.
  \item Catatan peer review lengkap.
\end{checklist}

P2-A menjadi bagian Praktikum 2 (12\%) yang dilaporkan pada Bab X.
